\documentclass[10pt,a4paper]{amsart}
\usepackage{amsmath,amssymb,amsthm,graphicx}
% Retained minimal preamble for the Stacks Project chapter more-algebra.tex (More on Algebra).
% The chapter file pinned-original-source.tex carries only \input{preamble} on its first line, so no
% preamble text lives inside the chapter itself.  The Stacks Project uses one single shared preamble for
% every chapter; the snapshot stages it identically as sources/stacks-chow/preamble.tex and
% sources/stacks-etale-cohomology/preamble.tex.  The block below is copied mechanically, line for line,
% from sources/stacks-chow/preamble.tex lines 14-18 (the Xy-pic setup) and 154-217 (the native theorem
% declarations together with the mathematical macros).  Nothing was added, removed or redefined; the
% wrapper class and the amsmath/amssymb/amsthm/graphicx base packages are supplied by the assembler.
% Chapter: more-algebra.tex -- More on Algebra.
% For commutative diagrams we use Xy-pic
\usepackage[all]{xy}

% We use 2cell for 2-commutative diagrams.
\xyoption{2cell}

% Theorem environments.
%
\theoremstyle{plain}
\newtheorem{theorem}[subsection]{Theorem}
\newtheorem{proposition}[subsection]{Proposition}
\newtheorem{lemma}[subsection]{Lemma}

\theoremstyle{definition}
\newtheorem{definition}[subsection]{Definition}
\newtheorem{example}[subsection]{Example}
\newtheorem{exercise}[subsection]{Exercise}
\newtheorem{situation}[subsection]{Situation}

\theoremstyle{remark}
\newtheorem{remark}[subsection]{Remark}
\newtheorem{remarks}[subsection]{Remarks}

\numberwithin{equation}{subsection}

% Macros
%
\def\lim{\mathop{\mathrm{lim}}\nolimits}
\def\colim{\mathop{\mathrm{colim}}\nolimits}
\def\Spec{\mathop{\mathrm{Spec}}}
\def\Hom{\mathop{\mathrm{Hom}}\nolimits}
\def\Ext{\mathop{\mathrm{Ext}}\nolimits}
\def\SheafHom{\mathop{\mathcal{H}\!\mathit{om}}\nolimits}
\def\SheafExt{\mathop{\mathcal{E}\!\mathit{xt}}\nolimits}
\def\Sch{\mathit{Sch}}
\def\Mor{\mathop{\mathrm{Mor}}\nolimits}
\def\Ob{\mathop{\mathrm{Ob}}\nolimits}
\def\Sh{\mathop{\mathit{Sh}}\nolimits}
\def\NL{\mathop{N\!L}\nolimits}
\def\CH{\mathop{\mathrm{CH}}\nolimits}
\def\proetale{{pro\text{-}\acute{e}tale}}
\def\etale{{\acute{e}tale}}
\def\QCoh{\mathit{QCoh}}
\def\Ker{\mathop{\mathrm{Ker}}}
\def\Im{\mathop{\mathrm{Im}}}
\def\Coker{\mathop{\mathrm{Coker}}}
\def\Coim{\mathop{\mathrm{Coim}}}

% Boxtimes
%
\DeclareMathSymbol{\boxtimes}{\mathbin}{AMSa}{"02}

%
% Macros for moduli stacks/spaces
%
\def\QCohstack{\mathcal{QC}\!\mathit{oh}}
\def\Cohstack{\mathcal{C}\!\mathit{oh}}
\def\Spacesstack{\mathcal{S}\!\mathit{paces}}
\def\Quotfunctor{\mathrm{Quot}}
\def\Hilbfunctor{\mathrm{Hilb}}
\def\Curvesstack{\mathcal{C}\!\mathit{urves}}
\def\Polarizedstack{\mathcal{P}\!\mathit{olarized}}
\def\Complexesstack{\mathcal{C}\!\mathit{omplexes}}
% \Pic is the operator that assigns to X its picard group, usage \Pic(X)
% \Picardstack_{X/B} denotes the Picard stack of X over B
% \Picardfunctor_{X/B} denotes the Picard functor of X over B
\def\Pic{\mathop{\mathrm{Pic}}\nolimits}
\def\Picardstack{\mathcal{P}\!\mathit{ic}}
\def\Picardfunctor{\mathrm{Pic}}
\def\Deformationcategory{\mathcal{D}\!\mathit{ef}}
\usepackage{cleveref}
\crefname{originallemma}{original lemma}{original lemmas}
\Crefname{originallemma}{Original lemma}{Original lemmas}
\makeatletter
\newlabel{more-algebra-lemma-injective-amplitude}{{0A5T}{source}{}{sourcepointer-1}{}}
\newlabel{more-algebra-lemma-injective-amplitude@cref}{{[originallemma][0][]0A5T}{[1][0][]source}{}{}{}}
\makeatother
\crefname{originallemma}{original lemma}{original lemmas}
\Crefname{originallemma}{Original lemma}{Original lemmas}
\makeatletter
\newlabel{derived-lemma-subcategory-right-resolution}{{05T6}{source}{}{sourcepointer-2}{}}
\newlabel{derived-lemma-subcategory-right-resolution@cref}{{[originallemma][0][]05T6}{[1][0][]source}{}{}{}}
\makeatother
\crefname{originallemma}{original lemma}{original lemmas}
\Crefname{originallemma}{Original lemma}{Original lemmas}
\makeatletter
\newlabel{lemma-characterize-injective-bis}{{0AVF}{source}{}{sourcepointer-3}{}}
\newlabel{lemma-characterize-injective-bis@cref}{{[originallemma][0][]0AVF}{[1][0][]source}{}{}{}}
\makeatother
\title{Three equivalent descriptions of bounded injective amplitude}
\author{The Stacks Project Authors}
\date{stacks-project:more-algebra.tex:local-snapshot-20261012; GNU FDL 1.2 or later}
\begin{document}
\maketitle
\setcounter{subsection}{1}
\noindent\textit{Editorial note.} All delivered mathematical text is the author's own, copied from the pinned Stacks Project chapter file more-algebra.tex; no statement or proof step was rewritten, reordered or removed.\par\medskip
\noindent\textit{Editorial note.} The chapter file pinned-original-source.tex carries only \textbackslash{}input\{preamble\} on its first line, so no preamble text lives inside the chapter. The retained file sources/stacks-more-algebra/retained-preamble.tex is a mechanical, line-for-line copy of the project's single shared preamble as staged in the snapshot (Xy-pic setup plus theorem declarations and mathematical macros). Nothing was added, removed or redefined, and the excerpt serves only as preamble, never as a span coordinate.\par\medskip
\noindent\textit{Editorial note.} Both counted spans start at the \textbackslash{}begin line of their environment, so the statement and the author's proof are delivered whole; surrounding prose of the section is deliberately outside the counted spans and is carried as standard prerequisites instead.\par\medskip
\noindent\textit{Editorial note.} The citation fields (source\_id, source\_version, authors, source\_url, license.name, source\_title) are identical to those of every other card this lane has produced from this chapter, so the batch selector treats them as one source.\par\medskip
\noindent\textit{Allowed standard prerequisites (editorial).}\begin{itemize}
\item Injective amplitude of a complex, defined through the existence of a bounded complex of injectives quasi-isomorphic to it. defined earlier in the same section, source lines 17480-17528, before the counted statement. Supplies the property appearing in the first condition.
\item Computation of Ext groups by a complex of injectives, and the exactness of the functor Hom from the ring. used at proof lines 17542-17560. Gives the vanishing statements in the second and third conditions.
\item Existence of a right resolution of a bounded below complex by injective modules. the author cites Derived Categories, Lemma derived-lemma-subcategory-right-resolution at proof lines 17570-17572. Produces the injective resolution used to compare with the injective amplitude condition.
\end{itemize}\medskip
\ifdefined\hypertarget\hypertarget{sourcepointer-1}{}\fi
\noindent\textit{Source reference.} Original lemma 0A5T: focused-seeds/AGP/sources/stacks-snapshot-20261011/tags.txt, source lines 12821--12821. The lemma's own canonical tag, quoted so that readers can resolve the label inside the delivered text; it points at the counted statement itself and carries no proof content.\par\medskip
\ifdefined\hypertarget\hypertarget{sourcepointer-2}{}\fi
\noindent\textit{Source reference.} Original lemma 05T6: focused-seeds/AGP/sources/stacks-snapshot-20261011/tags.txt, source lines 7470--7470. Navigation pointer to a statement the counted proof refers to but does not contain; nothing load-bearing is delegated to it.\par\medskip
\ifdefined\hypertarget\hypertarget{sourcepointer-3}{}\fi
\noindent\textit{Source reference.} Original lemma 0AVF: focused-seeds/AGP/sources/stacks-snapshot-20261011/tags.txt, source lines 13683--13683. Navigation pointer to a statement the counted proof refers to but does not contain; nothing load-bearing is delegated to it.\par\medskip
\begin{lemma}
\label{lemma-injective-amplitude}
Let $R$ be a ring. Let $K$ be an object of $D(R)$. Let $a, b \in \mathbf{Z}$.
The following are equivalent
\begin{enumerate}
\item $K$ has injective-amplitude in $[a, b]$,
\item $\Ext^i_R(N, K) = 0$ for all $R$-modules $N$ and all
$i \not \in [a, b]$,
\item $\Ext^i(R/I, K) = 0$ for all ideals $I \subset R$ and
all $i \not \in [a, b]$.
\end{enumerate}
\end{lemma}

\begin{proof}
Assume (1). We may assume $K$ is the complex
$$
\ldots \to 0 \to I^a \to I^{a + 1} \to \ldots \to
I^{b - 1} \to I^b \to 0 \to \ldots
$$
where $I^i$ is an injective $R$-module for all $i \in \mathbf{Z}$.
In this case we can compute the ext groups by the complex
$$
\ldots \to 0 \to \Hom_R(N, I^a) \to \ldots \to
\Hom_R(N, I^b) \to 0 \to \ldots
$$
and we obtain (2). It is clear that (2) implies (3).

\medskip\noindent
Assume (3) holds. Choose a nonzero map $R \to H^n(K)$. Since $\Hom_R(R, -)$
is an exact functor, we see that
$\Ext^n_R(R, K) = \Hom_R(R, H^n(K)) = H^n(K)$.
We conclude that $H^n(K)$ is zero for $n \not \in [a, b]$.
In particular, $K$ is bounded below and we can choose a quasi-isomorphism
$$
K \to I^\bullet
$$
with $I^i$ injective for all $i \in \mathbf{Z}$ and
$I^i = 0$ for $i < a$. See Derived Categories, Lemma
\ref{derived-lemma-subcategory-right-resolution}.
Let $J = \Ker(I^b \to I^{b + 1})$. Then $K$ is quasi-isomorphic
to the complex
$$
\ldots \to 0 \to I^a \to \ldots \to I^{b - 1} \to J \to 0 \to \ldots
$$
Denote $K' = (I^a \to \ldots \to I^{b - 1})$ the corresponding object of
$D(R)$. We obtain a distinguished triangle
$$
J[-b] \to K \to K' \to J[1 - b]
$$
in $D(R)$. Thus for every ideal $I \subset R$ an exact sequence
$$
\Ext^b(R/I, K') \to \Ext^1(R/I, J) \to \Ext^{1 + b}(R/I, K)
$$
By assumption the term on the right vanishes. By the implication
(1) $\Rightarrow$ (2) the term on the left vanishes. Thus $J$
is an injective $R$-module by
Lemma \ref{lemma-characterize-injective-bis}.
\end{proof}

\end{document}
